\documentclass[11pt,a4paper]{article}
\usepackage[margin=0.85in,top=1in,bottom=1in]{geometry}
\usepackage{amsmath,amssymb,amsfonts}
\usepackage{graphicx}
\usepackage{float}
\usepackage{enumitem}
\usepackage{microtype}
\usepackage{caption}
\usepackage[hidelinks]{hyperref}
\usepackage[most,breakable]{tcolorbox}
\tcbuselibrary{breakable,skins}
\usepackage[utf8]{inputenc}
\usepackage[T1]{fontenc}
\usepackage{lmodern}
\captionsetup{font=small,labelfont=bf,skip=4pt}
\setlist{leftmargin=1.5em,topsep=2pt}

%% Breakable card — prevents content cutoff across pages
\newtcolorbox{solcard}[3]{
  breakable, enhanced,
  colback=white, colframe=gray!50, arc=2pt,
  title={\textbf{#1}\hfill\textsf{\small #3\,pts}},
  fonttitle=\normalsize\sffamily\bfseries,
  before upper={\small\textit{#2}\par\smallskip\hrule height 0.4pt\smallskip},
  top=4pt, bottom=6pt, left=7pt, right=7pt,
  parbox=false,
}

\title{\textbf{Sun spots --- Solution}}
\author{\normalsize X22 (2022) — English translation}
\date{}
\begin{document}\maketitle\vspace{-0.5em}
\begin{solcard}{A1}{You are presented with a table of experimental data $\omega(\varphi)$. The data was obtained based on the Doppler effect, which affects the frequency of electromagnetic radiation from the Sun recorded on Earth. Define parameters $a$ and $b$.}{2.00}

Let's build a graph $\omega(\sin^2 \varphi)$. We draw an approximating line along the initial section, from which we find the parameters $a$ and $b$.

\par\smallskip\noindent\textbf{Answer:} $a = (2.99 \pm 0.2) \cdot 10^{-6}~1/\text{s}$
\begin{figure}[H]
  \centering
  \includegraphics[width=0.72\textwidth]{PE_S_EN_assets/fig1.jpg}
\end{figure}
\end{solcard}

\begin{solcard}{A2}{Find the limits of applicability of the formula given above.}{0.30}

From the graph we find the angle $\varphi_{\text{max}}$ after which the points significantly deviate from the straight line.

\par\smallskip\noindent\textbf{Answer:} $\varphi_{\text{max}} = 50^{\circ}$
\end{solcard}

\begin{solcard}{B1}{Obtain all components of the magnetic field on the surface of the Sun $B_{0,r}(\phi)$, $B_{0,\phi}(\phi)$, $B_{0,\lambda} (\phi)$. Express your answers using $\varphi$, $\mathfrak{M}$, $R_0$ and $\mu_0$.}{0.30}

Let us write the field of the magnetic dipole at latitude $\varphi$: \[ \vec{B}(\varphi) = \frac{\mu_0}{4\pi} \frac{\mathfrak{M}}{R_0^3} (3\vec{n} \sin \varphi - \vec{z} ) \] Projecting this onto the axes corresponding to $r$, $\varphi$ and $\lambda$ we obtain:

\par\smallskip\noindent\textbf{Answer:} \[B_{0,r} (\varphi) = \dfrac{\mu_0}{4\pi} \dfrac{\mathfrak{M}}{R_0^3} 2 \sin \varphi\]
\end{solcard}

\begin{solcard}{B2}{Obtain the magnitude of the magnetic field $B_e$ at the equator inside the Sun near its surface. Express your answer using $\mathfrak{M}$, $\mu_0$, $R_0$, $\xi$.}{0.60}

Let's choose a surface stretched over magnetic lines. Then the magnetic field flux through the equatorial section inside the Sun is equal to the flux through one of the hemispheres.


Thus $2\pi R_0 H_0 B_e = \Phi_0$ means $B_e = \dfrac{\Phi_0}{2\pi \xi R_0^2}=53~\text{mT}$. On the other hand, we can calculate the flux $\Phi_0$ by integrating the dipole field: \[ \Phi_0 = \int\limits_0^{\pi/2} B_{0,r} 2 \pi R_0^2 \cos \varphi d\varphi = \frac{\mu_0 \mathfrak{M}}{R_0} \int\limits_0^{\pi/2} \sin \varphi \cos \varphi d\varphi = \frac{\mu_0 \mathfrak{M}}{2R_0}. \] Substitute this into the previous formula and get the following:

\par\smallskip\noindent\textbf{Answer:} \[ B_e = \frac{\mu_0 \mathfrak{M}}{4\pi\xi R_0^3}. \]
\begin{figure}[H]
  \centering
  \includegraphics[width=0.72\textwidth]{PE_S_EN_assets/fig2.jpg}
\end{figure}
\end{solcard}

\begin{solcard}{B3}{Find the numerical value of $B_e$.}{0.10}
\par\smallskip\noindent\textbf{Answer:} \[B_e = \dfrac{\Phi_0}{2\pi \xi R_0^2}=53~\text{mT}\]
\end{solcard}

\begin{solcard}{B4}{Obtain the magnitude of the magnetic field $B_\varphi(\varphi)$ inside the Sun at its surface as a function of latitude $\varphi$. Express your answer using $\varphi$, $B_e$.}{0.60}

Now we will also select a surface stretched over magnetic lines. Only now it will extend from the equator to the desired latitude $\varphi$.


Then to calculate $B_\varphi$ you need to find the flux $\Phi(\varphi)$ of the magnetic field of the dipole through the part of latitude from $0$ to $\varphi$. This is the integral already familiar to us: \[\Phi(\varphi) (\varphi)= \int\limits_0^{\varphi} B_{0,r} 2 \pi R_0^2 \cos \varphi d\varphi = \frac{\mu_0 \mathfrak{M}}{2R_0} \sin^2 \varphi. \] The flux through the surface is equal to 0: \[ \Phi(\varphi) + 2\pi R_0 \cos \varphi \, H_0 \cos^2 \varphi\, B_\varphi(\varphi) = \Phi_0, \] then: \[ 2\pi\xi R_0^2 \cos^3 \varphi \,B_\varphi (\varphi) = \Phi_0 \cos^2 \varphi. \] So

\par\smallskip\noindent\textbf{Answer:} \[ B_\varphi = \frac{B_e}{\cos \varphi}.\]
\begin{figure}[H]
  \centering
  \includegraphics[width=0.72\textwidth]{PE_S_EN_assets/fig3.jpg}
\end{figure}
\end{solcard}

\begin{solcard}{C1}{Precisely calculate 30 pairs of points $(\varphi, \lambda)$ through which the magnetic line piercing points $\varphi=\pm 50^\circ$, $\lambda=0$ will pass at time $t=150~\text{days}$.}{1.70}

Let the magnetic line pass through the points $(\pm\varphi_0,\lambda_0)$, where $\varphi_0=50^\circ$, $\lambda_0=0^\circ$. Then its equation in angular coordinates is given as follows: \[ \lambda = \omega(\varphi)t = (a - b\sin^2 \varphi)t \]

\end{solcard}

\begin{solcard}{C2}{Using this data, plot on your answer sheet what this line will look like for an earthly observer.}{1.00}

Points with angular coordinates in the plane have the following coordinates: \[ \begin{cases} x = R \cdot \sin \lambda\, \cos \varphi\\ y = R \cdot \sin \varphi \end{cases} \] Let’s draw a line through the recalculated points:

\begin{figure}[H]
  \centering
  \includegraphics[width=0.72\textwidth]{PE_S_EN_assets/fig4.jpg}
\end{figure}
\end{solcard}

\begin{solcard}{C3}{Find the angle $\psi (\varphi,t)$ that the magnetic lines make with the meridian. Express your answer using $\varphi$, $t$, $a$, $b$.}{0.30}

From the picture it is clear that \[ \tan \psi = \frac{d \lambda}{d \varphi} \cos \varphi = 2bt \sin \varphi \cos^2 \varphi, \] is from

\par\smallskip\noindent\textbf{Answer:} \[ \psi = \arctan \left( 2bt \sin \varphi \, \cos^2 \varphi \right). \]
\begin{figure}[H]
  \centering
  \includegraphics[width=0.72\textwidth]{PE_S_EN_assets/fig5.jpg}
\end{figure}
\end{solcard}

\begin{solcard}{C4}{Let the magnetic lines inside the Sun near its surface make an angle $\psi$ with the meridian. Find the magnitude of the magnetic field $B$ at this point. Express your answer using $\psi$, $\varphi$, $B_e$.}{0.20}

The magnetic field flux through the surfaces discussed in Part B should not change, so $B_\varphi=\text{const}$. Then \[ B = \frac{B_\varphi}{\cos \psi} \] and therefore

\par\smallskip\noindent\textbf{Answer:} \[\frac{B_e}{\cos \varphi \, \cos \psi} \]
\end{solcard}

\begin{solcard}{D1}{Using the $\sin \psi \approx 1 $ approximation, find $B_\text{crit}$.}{2.70}

In the approximation $\sin \psi \approx 1$ it is true that $\tan \psi = 1/\cos \psi$, therefore: \[ B = \frac{B_e \tan \psi}{\cos \varphi} = 2B_e bt \sin \varphi, \] means the type of curve, $\varphi(t)$, on which spots appear on the Sun is described by the following expression: \[ B_\text{crit} = 2B_e bt \sin \varphi, \] i.e. its linearization is as follows: \[ t = \frac{A}{\sin 2\varphi} , \] where $A = \frac{B_{\text{crit}}}{B_eb}$ and $t$ are counted from the beginning of the 11-year solar cycle (one “butterfly” on the Maunder diagram). By constructing the envelope on the Maunder diagram, we obtain the dependence $t(\sin 2\varphi)$ for critical points (appearance of spots). From it we find $A$, and then $B_{crit}$.

\par\smallskip\noindent\textbf{Answer:} \[ B_{crit} = (15 \pm 5) \, \mathrm{mT} \]
\begin{figure}[H]
  \centering
  \includegraphics[width=0.72\textwidth]{PE_S_EN_assets/fig6.jpg}
\end{figure}
\end{solcard}

\begin{solcard}{D2}{Determine the period of solar activity $T_0$, i.e. the time after which the Sun's magnetic field returns to its original state.}{0.20}

We find the period as the time between two points with the same “phase” divided by the number of cycles: $T = 11~\text{year}$. Next, this number must be increased by 2 times, since after one cycle the field changes direction.

\par\smallskip\noindent\textbf{Answer:} \[ T_0 = 22 \, year \]
\end{solcard}

\end{document}